\section{Related Work}\label{sec:related}

Our work sits at the intersection of four recent threads in cache-eviction
research. \textbf{FIFO-based eviction:} Yang et al.\ (\cite{yang2023fifo})
proposed S3-FIFO as a three-queue FIFO replacement for LRU and demonstrated
that, with a small ``small-queue'' admitting one-hit wonders and promoting
survivors to the main queue, pure FIFO matches or beats LRU on Twemcache and
Memcached workloads. Zhang et al.\ (\cite{zhang2024sieve}) pushed this further
with SIEVE, a single-queue variant whose visited-bit scanning delivers
scan-resistance while remaining simpler than S3-FIFO's queue machinery.
\textbf{Admission filtering:} Einziger et al.\ (\cite{einziger2017tinylfu})
introduced TinyLFU, a Count-Min-Sketch-backed admission filter that forces new
candidates to out-frequency the eviction victim. Caffeine
(\cite{caffeine2024}) is the reference production implementation, with a 1\%
window-LRU that gives one-hit wonders a chance to prove recurrence.
\textbf{Scalable MRC construction:} Waldspurger et al.\
(\cite{waldspurger2015shards}) introduced SHARDS, a hash-based sampling
technique that makes miss-ratio-curve construction tractable at
million-request scale by reducing reuse-distance accounting to a sampled
subset. Our contribution is orthogonal to each of these: rather than proposing
a new policy, we compare six policies (LRU, FIFO, CLOCK, S3-FIFO, SIEVE,
W-TinyLFU) head-to-head on two real public-records API workloads with 5-seed
Welch's-t-backed confidence bands, and use SHARDS self-convergence to
validate our MRC construction at 1M-access scale where no exact oracle is
feasible.
